\documentclass[10pt,a4paper]{amsart}
\usepackage{amsmath,amssymb,amsthm,graphicx}
\usepackage{verbatim}
\newenvironment{reference}{\comment}{\endcomment}
%\newenvironment{reference}{}{}
\newenvironment{slogan}{\comment}{\endcomment}
\newenvironment{history}{\comment}{\endcomment}

% For commutative diagrams we use Xy-pic
\usepackage[all]{xy}

% We use 2cell for 2-commutative diagrams.
\xyoption{2cell}
% Theorem environments.
%
\theoremstyle{plain}
\newtheorem{theorem}[subsection]{Theorem}
\newtheorem{proposition}[subsection]{Proposition}
\newtheorem{lemma}[subsection]{Lemma}

\theoremstyle{definition}
\newtheorem{definition}[subsection]{Definition}
\newtheorem{example}[subsection]{Example}
\newtheorem{exercise}[subsection]{Exercise}
\newtheorem{situation}[subsection]{Situation}

\theoremstyle{remark}
\newtheorem{remark}[subsection]{Remark}
\newtheorem{remarks}[subsection]{Remarks}

\numberwithin{equation}{subsection}

% Macros
%
\def\lim{\mathop{\mathrm{lim}}\nolimits}
\def\colim{\mathop{\mathrm{colim}}\nolimits}
\def\Spec{\mathop{\mathrm{Spec}}}
\def\Hom{\mathop{\mathrm{Hom}}\nolimits}
\def\Ext{\mathop{\mathrm{Ext}}\nolimits}
\def\SheafHom{\mathop{\mathcal{H}\!\mathit{om}}\nolimits}
\def\SheafExt{\mathop{\mathcal{E}\!\mathit{xt}}\nolimits}
\def\Sch{\mathit{Sch}}
\def\Mor{\mathop{\mathrm{Mor}}\nolimits}
\def\Ob{\mathop{\mathrm{Ob}}\nolimits}
\def\Sh{\mathop{\mathit{Sh}}\nolimits}
\def\NL{\mathop{N\!L}\nolimits}
\def\CH{\mathop{\mathrm{CH}}\nolimits}
\def\proetale{{pro\text{-}\acute{e}tale}}
\def\etale{{\acute{e}tale}}
\def\QCoh{\mathit{QCoh}}
\def\Ker{\mathop{\mathrm{Ker}}}
\def\Im{\mathop{\mathrm{Im}}}
\def\Coker{\mathop{\mathrm{Coker}}}
\def\Coim{\mathop{\mathrm{Coim}}}

% Boxtimes
%
\DeclareMathSymbol{\boxtimes}{\mathbin}{AMSa}{"02}

%
% Macros for moduli stacks/spaces
%
\def\QCohstack{\mathcal{QC}\!\mathit{oh}}
\def\Cohstack{\mathcal{C}\!\mathit{oh}}
\def\Spacesstack{\mathcal{S}\!\mathit{paces}}
\def\Quotfunctor{\mathrm{Quot}}
\def\Hilbfunctor{\mathrm{Hilb}}
\def\Curvesstack{\mathcal{C}\!\mathit{urves}}
\def\Polarizedstack{\mathcal{P}\!\mathit{olarized}}
\def\Complexesstack{\mathcal{C}\!\mathit{omplexes}}
% \Pic is the operator that assigns to X its picard group, usage \Pic(X)
% \Picardstack_{X/B} denotes the Picard stack of X over B
% \Picardfunctor_{X/B} denotes the Picard functor of X over B
\def\Pic{\mathop{\mathrm{Pic}}\nolimits}
\def\Picardstack{\mathcal{P}\!\mathit{ic}}
\def\Picardfunctor{\mathrm{Pic}}
\def\Deformationcategory{\mathcal{D}\!\mathit{ef}}
\usepackage{cleveref}
\crefname{originallemma}{original lemma}{original lemmas}
\Crefname{originallemma}{Original lemma}{Original lemmas}
\makeatletter
\newlabel{algebra-lemma-completion-generalities}{{0315}{source}{}{sourcepointer-1}{}}
\newlabel{algebra-lemma-completion-generalities@cref}{{[originallemma][0][]0315}{[1][0][]source}{}{}{}}
\makeatother
\crefname{originallemma}{original lemma}{original lemmas}
\Crefname{originallemma}{Original lemma}{Original lemmas}
\makeatletter
\newlabel{algebra-lemma-quotient-complete}{{031A}{source}{}{sourcepointer-2}{}}
\newlabel{algebra-lemma-quotient-complete@cref}{{[originallemma][0][]031A}{[1][0][]source}{}{}{}}
\makeatother
\crefname{originalreference}{original reference}{original references}
\Crefname{originalreference}{Original reference}{Original references}
\makeatletter
\newlabel{lemma-noncomplete-quotient}{{05JE}{source}{}{sourcepointer-3}{}}
\newlabel{lemma-noncomplete-quotient@cref}{{[originalreference][0][]05JE}{[1][0][]source}{}{}{}}
\makeatother
\title{Completion is neither left nor right exact in general}
\author{The Stacks Project Authors}
\date{stacks-project:examples.tex:local-snapshot-20261012; GNU FDL 1.2 or later}
\begin{document}
\maketitle
\noindent\textit{Editorial note.} All delivered mathematical text is the author's own, copied from the pinned Stacks Project chapter file examples.tex; no statement or proof step was rewritten, reordered or removed.\par\medskip
\noindent\textit{Editorial note.} The counted proof is the author's complete argument: The statement's slogan environment is kept verbatim. The second paragraph's load-bearing ingredient (noncompleteness of R/J) is the delivered card L48305 in the same unit pool.\par\medskip
\noindent\textit{Editorial note.} The label lemma-completion-not-exact names the counted statement, and the proof environment immediately following it in the source is delivered verbatim as its proof.\par\medskip
\noindent\textit{Editorial note.} The counted statement is carried by the author's own lemma environment in the chapter source.\par\medskip
\noindent\textit{Allowed standard prerequisites (editorial).}\begin{itemize}
\item Context span 746-771, first paragraph: R = k[t], I = (t), P = K = direct sum over N of R, M = direct sum of R/(t\textasciicircum{}n); the short exact sequence 0 -> K -> P -> M -> 0 with first map multiplication by t\textasciicircum{}n on the nth summand; after completion, xi = (t\textasciicircum{}2, t\textasciicircum{}3, t\textasciicircum{}4, ...) in P\textasciicircum{}wedge maps to zero in M\textasciicircum{}wedge but is not in the image of K\textasciicircum{}wedge -> P\textasciicircum{}wedge since that would be the image of (t, t, t, ...) which is not an element of K\textasciicircum{}wedge
\item Context, second paragraph: in the situation of Lemma noncomplete-quotient (delivered card L48305) the short exact sequence 0 -> J -> R -> R/J -> 0 does not remain exact after I-adic completion: with f a generator of J the map R -> J\textasciicircum{}wedge is surjective with image of J\textasciicircum{}wedge -> R equal to (f) = J, but R/J noncomplete means the kernel of the surjection R -> (R/J)\textasciicircum{}wedge is strictly bigger than J (Algebra, completion-generalities and quotient-complete); by the same token R -> R -> R/(f) -> 0 does not remain exact on completion
\end{itemize}\medskip
\ifdefined\hypertarget\hypertarget{sourcepointer-1}{}\fi
\noindent\textit{Source reference.} Original lemma 0315: sources/stacks-chow/tags.txt, source lines 3972--3972. Prerequisite pointer invoked at 808; it names a standard fact of a sibling chapter that the counted argument uses, and every step of the counted proof built on it is retained verbatim in the delivered body.\par\medskip
\ifdefined\hypertarget\hypertarget{sourcepointer-2}{}\fi
\noindent\textit{Source reference.} Original lemma 031A: sources/stacks-chow/tags.txt, source lines 3977--3977. Prerequisite pointer invoked at 809; it names a standard fact of a sibling chapter that the counted argument uses, and every step of the counted proof built on it is retained verbatim in the delivered body.\par\medskip
\ifdefined\hypertarget\hypertarget{sourcepointer-3}{}\fi
\noindent\textit{Source reference.} Original reference 05JE: sources/stacks-examples/pinned-original-source.tex, source lines 730--730. Prerequisite pointer invoked at 800; it names a standard fact of the same chapter that the counted argument uses, and every step of the counted proof built on it is retained verbatim in the delivered body.\par\medskip
A quick example is the following. Suppose that $R = k[t]$. Denote $M^\wedge$
the completion of an $R$-module with respect to $I = (t)$. Let
$P = K = \bigoplus_{n \in \mathbf{N}} R$ and
$M = \bigoplus_{n \in \mathbf{N}} R/(t^n)$. Then there is a short exact
sequence $0 \to K \to P \to M \to 0$ where the first map is given by
multiplication by $t^n$ on the $n$th summand. We claim that
$0 \to K^\wedge \to P^\wedge \to M^\wedge \to 0$ is not exact in the middle.
Namely, $\xi = (t^2, t^3, t^4, \ldots) \in P^\wedge$ maps to zero in
$M^\wedge$ but is not in the image of $K^\wedge \to P^\wedge$, because
it would be the image of $(t, t, t, \ldots)$ which is not an element of
$K^\wedge$.

\medskip\noindent
A ``smaller'' example is the following. In the situation of
Lemma \ref{lemma-noncomplete-quotient}
the short exact sequence $0 \to J \to R \to R/J \to 0$ does not remain
exact after completion with respect to $I$.
Namely, if $f \in J$ is a generator, then
$f : R \to J$ is surjective, hence $R \to J^\wedge$ is surjective, hence
the image of $J^\wedge \to R$ is $(f) = J$ but the fact that
$R/J$ is noncomplete means that the kernel of the surjection
$R \to (R/J)^\wedge$ is strictly bigger than $J$, see
Algebra, Lemmas \ref{algebra-lemma-completion-generalities} and
\ref{algebra-lemma-quotient-complete}.
By the same token the sequence
$R \to R \to R/(f) \to 0$ does not remain exact on completion.

\begin{lemma}
\label{lemma-completion-not-exact}
\begin{slogan}
Completion is neither left nor right exact in general.
\end{slogan}
Completion is not an exact functor in general; it is not even
right exact in general. This holds even when $I$ is finitely
generated on the category of finitely presented modules.
\end{lemma}

\begin{proof}
See discussion above.
\end{proof}

\end{document}
